\begin{proof}[Proof of \cref{obj:thm:matched-minimax-frontier}]
\begin{enumerate}
\item Choose universal constants \(\underline c\) and \(C_0\) supplied by
\cref{obj:thm:full-experiment-lower,obj:thm:mixed-count-upper}, respectively, and set
\[
0<\underline c,\qquad 0<C_0,\qquad
\overline C:=\max\{\underline c,C_0\}.
\]
Thus
\[
0<\underline c\le\overline C<\infty,
\qquad C_0\le\overline C.
\]
Fix \(n,d,q\) satisfying the hypotheses of the theorem.
By \cref{obj:thm:full-experiment-lower},
\[
\underline c\,r_{n,d,q}\le\mathfrak R_{n,d,q}.
\]
% lean: CausalSmith.Stat.MarRareqLogfrontier.matched_minimax_frontier

\item We next construct an admissible decision kernel from the deterministic estimator in
\cref{obj:def:mixed-count-estimator}. For each observed sample \(o\), its defining finite average is
\[
\widehat\tau^{\mathrm{mix}}_{n,d,q}(o)
=
\sum_{\nu=0}^{n}
\Pr_\omega(K_0=\nu)\,
E_\omega\!\left[
\widetilde\tau(o;\omega)\mid K_0=\nu
\right].
\]
The omitted Poisson tail contributes zero by the estimator's definition.
Each conditional expectation is the uniform average over the \(3^\nu\) stream assignments.
Clipping, together with the zero-output branches, gives
\[
-1\le \widetilde\tau(o;\omega)\le1,
\qquad
-1\le
E_\omega\!\left[
\widetilde\tau(o;\omega)\mid K_0=\nu
\right]
\le1
\quad(0\le\nu\le n).
\]
The prefix weights satisfy
\[
\Pr_\omega(K_0=\nu)\ge0,\qquad
\sum_{\nu=0}^{n}\Pr_\omega(K_0=\nu)
=\Pr_\omega(K_0\le n)\le1.
\]
Consequently, at every observed sample,
\[
-1
\le-\Pr_\omega(K_0\le n)
\le\widehat\tau^{\mathrm{mix}}_{n,d,q}(o)
\le\Pr_\omega(K_0\le n)
\le1.
\]
The observed-sample space is finite with measurable singletons, so this deterministic map is measurable. Its Dirac kernel,
\[
\mathsf T_{\mathrm{mix}}(o,dh)
:=
\delta_{\widehat\tau^{\mathrm{mix}}_{n,d,q}(o)}(dh),
\]
therefore belongs to the decision class \(\mathcal T_n\) of
\cref{obj:def:minimax-risk}.
% lean: CausalSmith.Stat.MarRareqLogfrontier.mixed_count_estimator_regular

\item For every \(\mathsf T\in\mathcal T_n\) and \(P\in\mathcal M_{n,d,q}\), the joint squared-error risk is nonnegative:
\[
E_P\!\left[(\mathsf T-\tau(P))^2\right]
=
\int\!\int_{[-1,1]}
(h-\tau(P))^2\,\mathsf T(o,dh)\,
\mathcal L_P(\mathbf O_n)(do)
\ge0,
\]
because both integrals have nonnegative integrands. For the Dirac kernel just constructed,
\[
E_P\!\left[(\mathsf T_{\mathrm{mix}}-\tau(P))^2\right]
=
E_P\!\left[
\left(
\widehat\tau^{\mathrm{mix}}_{n,d,q}(\mathbf O_n)-\tau(P)
\right)^2
\right].
\]
Moreover,
\[
nq>0,\qquad
r_{n,d,q}
=
\min\left\{
1,\,
\frac1{nq}
+
\left[\frac{d}{nq\log(e+nq)}\right]^2
\right\}
\ge0.
\]
For each \(P\in\mathcal M_{n,d,q}\),
\cref{obj:thm:mixed-count-upper} now yields the uniform bound
\[
E_P\!\left[(\mathsf T_{\mathrm{mix}}-\tau(P))^2\right]
\le \mathfrak U_{n,d,q}
\le C_0\,r_{n,d,q}
\le\overline C\,r_{n,d,q}.
\]
The last inequality uses \(C_0\le\overline C\) and \(r_{n,d,q}\ge0\).
% lean: CausalSmith.Stat.MarRareqLogfrontier.mixed_count_upper

\item Nonnegativity of all the risks permits comparison of the minimax value with the worst-case risk of any admissible estimator. In particular,
\[
\mathfrak R_{n,d,q}
\le
\sup_{P\in\mathcal M_{n,d,q}}
E_P\!\left[(\mathsf T_{\mathrm{mix}}-\tau(P))^2\right].
\]
If \(\mathcal M_{n,d,q}\) is nonempty, the uniform bound from the preceding step gives
\[
\sup_{P\in\mathcal M_{n,d,q}}
E_P\!\left[(\mathsf T_{\mathrm{mix}}-\tau(P))^2\right]
\le\overline C\,r_{n,d,q}.
\]
If \(\mathcal M_{n,d,q}\) is empty, the real-supremum convention assigns its worst-case risk the value zero. Since \(\overline C\ge C_0>0\) and \(r_{n,d,q}\ge0\), in this case
\[
\sup_{P\in\mathcal M_{n,d,q}}
E_P\!\left[(\mathsf T_{\mathrm{mix}}-\tau(P))^2\right]
=0\le\overline C\,r_{n,d,q}.
\]
Thus, in either case,
\[
\underline c\,r_{n,d,q}
\le\mathfrak R_{n,d,q}
\le\overline C\,r_{n,d,q}.
\]
The constants were chosen independently of \(n,d,q\), as required.
% lean: CausalSmith.Stat.MarRareqLogfrontier.matched_minimax_frontier
\end{enumerate}
\end{proof}
